Знания, содержащиеся в химических книгах, монографиях и архивных выпусках журналов, сохраняют высокую научную и практическую ценность: в них зафиксированы методики синтеза, экспериментальные наблюдения, справочные сведения о веществах и их свойствах, а также контекст, необходимый для воспроизводимости результатов. При этом существенная часть такого корпуса литературы до сих пор существует в виде физических изданий или в виде скан-копий, которые удобны для чтения человеком, но плохо приспособлены для автоматизированного анализа и интеграции в современные информационные системы.

На фоне развития информационных технологий и интеллектуального поиска становится актуальной задача извлечения знаний из химических книг с последующей структуризацией и верификацией. Однако на практике оцифровка часто сводится к распознаванию текста (OCR) и формированию <<поискового PDF>>~\cite{refDavletov}, что позволяет выполнять лишь поверхностный текстовый поиск. В химической литературе значимая часть информации представлена не в виде текста, а в виде графических объектов: структурных формул, схем реакций. Если эти элементы не переводятся в машинно-обрабатываемые представления, то большая доля смысла документа остается вне цифрового контура и фактически недоступна для алгоритмов поиска, сопоставления и анализа~\cite{refKuchuganov}.

В данной работе рассматривается подход к комплексной оцифровке химической литературы, ориентированный не только на извлечение текстового слоя, но и на распознавание и структурирование графической химической информации, что позволяет разработать систему интеллектуального поиска по книгам .